%% ==================== 问题重述 模板 (v1.5.6, BZD 数模社 2026 规范) ====================
%%
%% 借鉴 BZD 数模社 第一章 模板: 1.1 研究背景 + 1.2 问题回顾 + 1.3 研究综述
%%
%% 写作原则 (BZD 4 色教学法):
%% - 黑色 = 文章主体, 直接套用
%% - 蓝色 = 操作步骤
%% - 红色 = 解释说明
%% - 绿色 = 常见误区
%% - 黄色 = AI 提示词
%% - 灰色 = 自查清单
%%
%% 1.1 研究背景: 1-3 段, 研究对象 + 现实情境 + 主要矛盾 + 研究价值
%% 1.2 问题回顾: 按题号重述每问, 保留关键数据/单位/约束
%% 1.3 研究综述: ~500 字, 3-4 段, 现有方法 + 不足 + 本文切入点
%%
%% 跑题用户必须:
%% 1. 删掉所有 % 注释 (整段 L1-L29)
%% 2. 替换【】占位符为自己的实际内容
%% 3. 删 下面的 \AbstractSlot{...} 行 (check 15 抓这个 inline 【】)
%% 4. 跑题后用 LaTeX 编译一次验证 (xelatex 论文.tex × 2)

\section{问题重述}

\subsection{研究背景}

\textbf{【研究对象】}是指【简要定义, 1 句】, 在【应用领域, 1 句】中具有【主要作用, 1 句】。随着【社会/经济/技术/政策背景, 1 句】的发展, \textbf{【研究对象】}呈现出【主要趋势, 1 句】。

与此同时, 其发展仍受到【主要矛盾或现实问题, 1 句】的影响, 因此有必要对【核心研究问题, 1 句】进行系统研究。

\subsection{问题回顾}

根据题目提供的数据、条件和任务要求, 需要解决以下问题:

（1） 分析【问题一研究对象, 1 句】, 确定【需要识别/评价/计算的内容, 1 句】。

（2） 根据【数据条件, 1 句】, 建立模型完成【问题二的预测/评价/优化任务, 1 句】。

（3） 考虑【新增条件, 1 句】, 研究【问题三的主要任务, 1 句】。

（4） 结合【指定情境或可用的前述结果, 1 句】, 完成【问题四的任务, 1 句】。

若题目还有其他问题, 按照相同方式继续准确表述。

\subsection{研究综述}

现有研究主要从【方向一】、【方向二】和【方向三】分析该问题。在【方向一】方面, 常用【模型或方法】研究【具体任务】; 在【方向二】方面, 主要采用【模型或方法】分析【具体问题】; 在【方向三】方面, 相关研究通常通过【模型或方法】评价【研究对象】。

现有方法为研究该问题提供了基础, 但在【现有不足 1】、【现有不足 2】和【现有不足 3】方面仍具有一定局限。基于题目条件和研究任务, 有必要进一步从【本文建模角度】开展分析。

% 跑题后必须删下面这一行 (check 15 占位符自检)
\textbf{【这里粘贴从 1.1 到 1.3 拼起来的不超过 2 页内容, 删掉这一行】}
